\documentclass[11pt]{article}
\usepackage[a4paper,margin=2.4cm]{geometry}
\usepackage{amsmath,amssymb}
\usepackage{booktabs}
\usepackage{xcolor}
\usepackage{enumitem}
\usepackage[hidelinks]{hyperref}
\usepackage{siunitx}

\newcommand{\phin}{\phi}
\newcommand{\neqcm}{\ensuremath{\mathrm{n_{eq}/cm^{2}}}}
\newcommand{\fC}{\ensuremath{\mathrm{fC}}}
\newcommand{\code}[1]{\texttt{#1}}

\title{\textbf{Q(V,\,$\phi$) parametrization --- method}}
\author{ETL LGAD charge vs bias and fluence}
\date{\today}

\begin{document}
\maketitle

\noindent
This note describes, step by step, the laws used to parametrize the charge
$Q$ collected by ETL LGADs as a function of the bias $V$ and the
neutron--equivalent fluence $\phi$. It documents the code in this repository
(\code{data\_io.py}, \code{model.py}, \code{plots.py}, \code{config.py}).

\section{Data and pre-processing}
Measurements live in \code{parametrization\_data.xlsx}, one \emph{sheet per
production} (e.g.\ \code{HPK}; \code{FBK} later). Inside a sheet each sensor scan
is a block of three columns $V$, $Q$, $\mathrm{Res}$; the block label carries the
sensor name and the fluence (e.g.\ \code{HPK6 S4-B - 1e15}). The reader extracts:
\begin{itemize}[nosep]
  \item \textbf{production} $=$ sheet name;
  \item \textbf{fluence} parsed from the label ($0$ if none is written);
  \item \textbf{setup} $=$ \code{TB} if the label mentions a test beam, else \code{Beta}.
\end{itemize}

\paragraph{Pooling.} For a given production and fluence, the $(V,Q)$ points of all
\code{Beta} sensors are \emph{pooled} into a single data set and fitted together;
test--beam (\code{TB}) points are shown only as a cross--check and are never
fitted. Within one production sensors of nominally equivalent generations may be
pooled across fluence (e.g.\ HPK6 and HPK2).

\paragraph{Uncertainties.} The raw data has no error bars, so constant values are
assigned (editable in \code{config.py}):
\begin{equation}
  \sigma_Q = \SI{1}{\fC}\ \text{(switchable to \SI{2}{\fC})},\qquad
  \sigma_V = \SI{2}{\volt},\qquad
  \sigma_\phi = 0.10\,\phi .
\end{equation}
A constant $\sigma_Q$ does not bias the central fit values; it only sets the
weighting and the reported parameter errors.

\section{Step 1 --- charge vs bias at fixed fluence}\label{sec:step1}
At each fluence the pooled $Q(V)$ points are fitted with
\begin{equation}
  \boxed{\,Q(V) = p_0 + p_1\,V + \exp\!\big(p_2 + p_3\,V\big)\,}
  \label{eq:qv}
\end{equation}
a linear term plus an exponential (gain--like) rise.

\paragraph{Analytic seeding (no manual tuning).} ROOT's default start values make
this fit unstable, which previously required hand--tuned parameter limits. Instead
the seed is computed from the data (\code{auto\_seed\_QV}):
\begin{enumerate}[nosep]
  \item sort by $V$; fit a straight line to the lowest ${\sim}40\%$ of the points
        to get the baseline $(p_0,p_1)$;
  \item form the residual $R(V)=Q-(p_0+p_1V)$ and, on the points with $R>0$,
        fit $\ln R = p_2 + p_3 V$ to get $(p_2,p_3)$.
\end{enumerate}
With this seed the fit converges on its own.

\paragraph{Physical bounds.} The only constraints are \emph{uniform} (the same for
every fit, not per point):
\begin{equation}
  p_1 \ge 0, \qquad p_3 \ge 0,
\end{equation}
i.e.\ the collected charge does not decrease with bias. The fit is performed in
ROOT (\code{TGraphErrors::Fit} of Eq.~\ref{eq:qv}); if ROOT fails to return a
valid result, a \code{scipy.optimize.curve\_fit} with the same seed is used as a
silent fallback.

The outputs are the four parameters $p_k(\phi_j)$ and their errors at each measured
fluence $\phi_j$, with the fit $\chi^2/\mathrm{ndf}$.

\section{Step 2 --- parameters vs fluence}\label{sec:step2}
Each parameter $p_k$ is then modelled as a function of $\phi$. Four laws are
available (\code{config.PARAM\_MODE}); let $\{(\phi_j,p_{k,j},\sigma_{k,j})\}$ be
the points from Step~1.

\paragraph{(i) Node--exact interpolation \code{interp} (default).}
Piecewise--linear through the measured points, with linear extrapolation beyond
the end nodes:
\begin{equation}
  p_k(\phi) = p_{k,j} + \frac{p_{k,j+1}-p_{k,j}}{\phi_{j+1}-\phi_j}\,(\phi-\phi_j),
  \qquad \phi\in[\phi_j,\phi_{j+1}].
\end{equation}
By construction it reproduces every measured fluence exactly, so the predicted
$Q(V,\phi)$ matches the Step--1 fit at each $\phi_j$.

\paragraph{(ii) Linear \code{lin}.} Weighted least squares,
\begin{equation}
  p_k(\phi) = m_k\,\phi + q_k, \qquad
  \text{weights } w_j = 1/\sigma_{k,j}.
\end{equation}

\paragraph{(iii) Exponential \code{exp}.}
\begin{equation}
  p_k(\phi) = a_k + \exp\!\big(b_k\,\phi + c_k\big),
\end{equation}
fitted in ROOT with an analytic seed ($a_k$ near the asymptote, $(b_k,c_k)$ from a
log fit of $p_k-a_k$).

\paragraph{(iv) Automatic \code{auto}.} If there are fewer than
\code{MIN\_PTS\_FOR\_EXP} $=4$ fluences, use \code{lin}. Otherwise fit both
\code{lin} and \code{exp} and keep the one with the smaller Akaike information
criterion,
\begin{equation}
  \mathrm{AIC} = \chi^2 + 2k,
\end{equation}
with $k$ the number of free parameters ($2$ for \code{lin}, $3$ for \code{exp}).
This is decided independently for each $p_k$.

\paragraph{Why \code{interp} is the default here.} The form of
Eq.~\ref{eq:qv} is partly degenerate: $p_0,p_1,p_2$ are strongly
anti--correlated, so while the \emph{curve} $Q(V)$ is well constrained the
\emph{individual} parameters are not and scatter from fluence to fluence. As a
result the measured $p_k(\phi)$ do not follow a smooth low--order law, and a global
\code{lin}/\code{exp} fit of $p_2,p_3$ --- to which $Q$ is exponentially
sensitive --- can mispredict the charge badly. The node--exact \code{interp}
avoids this by reproducing the data at every measured fluence; \code{auto}/\code{exp}
remain available when a smooth, extrapolable surface is preferred.

\section{Combined surface and its inverse}
Substituting the parameter laws into Eq.~\ref{eq:qv} gives the two--dimensional
surface
\begin{equation}
  \boxed{\,Q(V,\phi) = P_0(\phi) + P_1(\phi)\,V +
          \exp\!\big(P_2(\phi) + P_3(\phi)\,V\big)\,}
\end{equation}
with $P_k(\phi)$ one of the laws above. When every $P_k$ is \code{lin} or
\code{exp} this is exported as a single closed--form ROOT \code{TF2}
(\code{build\_QVphi\_TF2}); for \code{interp} the surface is evaluated
numerically (\code{Q\_of\_V\_phi}).

\paragraph{Bias for a target charge.} To get the bias needed to reach a charge
$Q^\star$ at a fluence $\phi$, the equation $Q(V,\phi)=Q^\star$ is inverted
numerically (\code{V\_for\_Q}): $Q(V,\phi)$ is scanned on a dense $V$ grid and the
\emph{highest} upward crossing of $Q^\star$ is returned (robust to any
non--monotonic behaviour of a degenerate fit); \code{NaN} if $Q^\star$ is never
reached.

\section{The analytic model: gain term with a physical floor}
\label{sec:global}
The four--parameter form of Eq.~\ref{eq:qv} is over--parametrized for this data:
the per--fluence fits give $p_0$ and $p_1$ with relative errors of several hundred
to several million percent (they are not determined), and $p_0,p_1,p_2$ are
strongly anti--correlated. Using independent laws for all four makes the surface
diverge (prediction RMS $\sim\!28\,\fC$ vs $\sim\!1.5\,\fC$ below). The data
constrain only the \emph{gain} term, so the model keeps just that:
\begin{equation}
  \boxed{\,Q(V,\phi) = \exp\!\big(p_2(\phi) + p_3(\phi)\,V\big)\,}
  \label{eq:global}
\end{equation}

\paragraph{Physical floor.} The collected charge never drops below the no--gain
(bulk MIP) charge $Q_{\mathrm{floor}} = \SI{0.55}{\fC}$ (\code{config.Q\_FLOOR}):
at very high fluence, whatever the voltage, the minimum is $Q_{\mathrm{floor}}$.
This is imposed as a \emph{floor} (a clamp), not as an additive constant, so it
does not inflate $Q$ in the operating range:
\begin{equation}
  Q_{\mathrm{phys}}(V,\phi) = \max\!\big(Q_{\mathrm{floor}},\,
      \exp(p_2(\phi)+p_3(\phi)V)\big).
\end{equation}
Because the gain exponent decreases with fluence, $\exp(\dots)\to0$ at high $\phi$
and the clamp delivers a flat $Q_{\mathrm{floor}}$ for every $V$ --- consistent
with the physics, and impossible with a $p_0$ (which would set the floor at $0\,$V
too) or a $p_1 V$ term (which would make the high--$\phi$ charge grow with $V$).

\paragraph{Step A (per fluence).} $p_2,p_3$ are obtained from a \emph{linear
regression of} $\ln Q$ \emph{vs} $V$ (the exact, bound--free, seed--free fit for an
exponential; weight $Q^2$ for the $\sigma_Q=$const error). This is stable where a
direct $\exp$ fit can diverge.

\paragraph{Step B (vs fluence).} Each of $p_2(\phi),p_3(\phi)$ gets an analytic
law (linear or exponential, chosen by AIC, Sec.~\ref{sec:step2}; \code{interp}
is never used here). For HPK both come out exponential. The bias for a target
charge $Q^\star$ is explicit, $V(\phi) = [\ln Q^\star - p_2(\phi)]/p_3(\phi)$.
This model is \code{model.GlobalSurface} / \code{model.fit\_global}; the
four--parameter per--fluence fits of Sec.~\ref{sec:step1} remain as diagnostics.

\section{Operating point with the SEB limit}
Each production has a single--event--burnout ceiling $V_{\mathrm{SEB}}$
(\code{config.V\_SEB}; HPK \SI{550}{\volt}, FBK \SI{600}{\volt}) above which it
cannot be operated. For a target charge $Q^\star$ (default
$\code{PLATEAU\_CHARGE}=\SI{15}{\fC}$):
\begin{align}
  \text{bias needed:}\quad & V(\phi) = \min\!\big(V_{\mathrm{SEB}},\,
      V_{\mathrm{for}\,Q^\star}(\phi)\big),\\
  \text{charge obtainable:}\quad & Q_{\mathrm{obt}}(\phi) =
      \min\!\big(Q^\star,\, Q(V_{\mathrm{SEB}},\phi)\big).
\end{align}
Below the ``SEB knee'' the target is reached within $V_{\mathrm{SEB}}$ so
$Q_{\mathrm{obt}}=Q^\star$ (flat); beyond it the device is stuck at
$V_{\mathrm{SEB}}$ and delivers the smaller charge $Q(V_{\mathrm{SEB}},\phi)$,
which falls with fluence towards $Q_{\mathrm{floor}}$. Here $Q(V_{\mathrm{SEB}},\phi)$
is taken from the global surface of Sec.~\ref{sec:global}, so the curve is a single
smooth trend and the charge converges to $Q_{\mathrm{floor}}=\SI{0.55}{\fC}$, not to
zero. These are the solid (charge, left axis) and dashed (bias, right axis) curves
of the operating--point figure.

\section{Caveats}
\begin{itemize}[nosep]
  \item Parameter degeneracy (above) makes the individual $p_k$ noisy; prefer the
        \emph{curves}/predictions over the individual parameter values.
  \item Extrapolation in $V$ beyond the measured range, and in $\phi$ beyond the
        last measured fluence, is model dependent and should be treated with care.
  \item With \code{interp} the surface is only as smooth as the data; for a smooth
        analytic law use \code{auto}/\code{exp}, accepting that it will not pass
        exactly through every measured fluence.
\end{itemize}

\end{document}
